\section{Experimental Design}
\label{sec:experimental_design}

Following the experimental software engineering methodology of Wohlin et al. \cite{wohlin_2012} and the GQM framework \cite{basili_1994}, our experimental goal is to \textbf{analyze} Chaos Mesh and LitmusChaos; \textbf{for the purpose of} comparative empirical evaluation; \textbf{with respect to} temporal actuation/teardown latency ($\Delta T_{\text{actuate}}, \Delta T_{\text{teardown}}$), cluster resource overhead ($\Delta_{\text{CPU}}, \Delta_{\text{RAM}}$), multimodal RCA diagnostic parity ($MTTDiag, Acc$), and native probe verification expressiveness; \textbf{from the point of view of} cloud-native architects and SREs; \textbf{in the context of} a Kubernetes microservice cluster instrumented with OpenTelemetry, Prometheus, Loki, and Tempo.

\subsection{Hypotheses Formulation}
Corresponding to our four Research Questions (RQ01--RQ04), we test four formal hypotheses ($\alpha = 0.05$):
\begin{itemize}[leftmargin=1.2em, itemsep=2pt, topsep=3pt]
    \item \textbf{Actuation \& Teardown Latency Hypothesis ($H_{1,1}$ — RQ01):} Chaos Mesh exhibits significantly lower fault actuation delay ($\Delta T_{\text{actuate}}^{\text{CM}} \allowbreak<\allowbreak \Delta T_{\text{actuate}}^{\text{LC}}$) and zero ephemeral runner pod churn compared to LitmusChaos due to resident DaemonSet execution versus two-stage runner pod scheduling.
    \item \textbf{Resource Overhead Hypothesis ($H_{1,2}$ — RQ02):} In standalone CRD mode (\texttt{litmus-core}), LitmusChaos consumes lower idle control-plane memory than Chaos Mesh's node-resident DaemonSet ($\Delta_{\text{RAM,idle}}^{\text{LC}} \allowbreak<\allowbreak \Delta_{\text{RAM,idle}}^{\text{CM}}$), whereas deploying LitmusChaos's full stateful control plane (\texttt{ChaosCenter}/MongoDB) or active runner Jobs shifts the resource trade-off toward transient pod-scheduling overhead.
    \item \textbf{Diagnostic Parity \& Telemetry Noise Hypothesis ($H_{1,3}$ — RQ03):} When referenced to actual fault activation ($T_{\text{effect}}$), equivalent fault profiles in Chaos Mesh and LitmusChaos yield statistically equivalent root-cause diagnostic accuracy ($Acc^{\text{CM}} \allowbreak\approx\allowbreak Acc^{\text{LC}}$) in the PANOPTES multimodal RCA engine, provided that transient \texttt{chaos-runner}/\texttt{go-runner} pods are excluded via namespace/label filtering.
    \item \textbf{Native Probe Expressiveness Hypothesis ($H_{1,4}$ — RQ04):} LitmusChaos's multi-modal native probes (\texttt{promProbe} + \texttt{k8sProbe} + \texttt{httpProbe}) detect a significantly higher proportion of process-invisible (\texttt{cpu-stress}) and lifecycle (\texttt{pod-kill}) perturbations autonomously than Chaos Mesh's HTTP-only \texttt{StatusCheck} CRD ($Acc_{\text{native}}^{\text{LC}} \allowbreak>\allowbreak Acc_{\text{native}}^{\text{CM}}$).
\end{itemize}

\begin{table*}[!tbp]
\centering
\scriptsize
\setlength{\tabcolsep}{3.5pt}
\renewcommand{\arraystretch}{1.15}
\caption{Equivalent Fault Campaign Mapping Between Chaos Mesh and LitmusChaos}
\label{tab:fault_equivalence}
\begin{tabularx}{\textwidth}{@{}>{\raggedright\arraybackslash}p{2.35cm} >{\raggedright\arraybackslash}p{2.45cm} >{\raggedright\arraybackslash}X >{\raggedright\arraybackslash}X@{}}
\toprule
\textbf{Fault Stratum ($n$)} & \textbf{Target Pool ($\mathcal{P}_{\text{target}}$)} & \textbf{Chaos Mesh v2.6.3 Specification} & \textbf{LitmusChaos v3.x Equivalent Specification} \\
\midrule
\textbf{1. Network Latency} ($n=10$) & \texttt{cart}, \texttt{currency}, \texttt{payment}, \texttt{shipping} & \texttt{kind: NetworkChaos}, \texttt{action: delay}, \texttt{latency: "2000ms"}, \texttt{jitter: "100ms"}, \texttt{direction: to}, \texttt{duration: "180s"} (via \texttt{chaos-daemon} \texttt{tc-netem}). & \texttt{ChaosExperiment: pod-network-latency}, \texttt{NETWORK\_\allowbreak LATENCY: "2000"}, \texttt{JITTER: "100"}, \texttt{TOTAL\_\allowbreak CHAOS\_\allowbreak DURATION: "180"} (via \texttt{go-runner} \texttt{tc-netem}). \\
\textbf{2. Pod Termination} ($n=9$) & \texttt{cart}, \texttt{currency}, \texttt{payment}, \texttt{shipping} & \texttt{kind: PodChaos}, \texttt{action: pod-kill}, \texttt{mode: one}, explicit pod-instance selector (\texttt{spec.\allowbreak selector.\allowbreak pods}). & \texttt{ChaosExperiment: pod-delete}, \texttt{FORCE: "true"}, \texttt{TOTAL\_\allowbreak CHAOS\_\allowbreak DURATION: "15"}, \texttt{TARGET\_\allowbreak PODS: <pod-name>}. \\
\textbf{3. CPU Saturation} ($n=11$) & \texttt{cart}, \texttt{currency}, \texttt{payment}, \texttt{shipping} & \texttt{kind: StressChaos}, \texttt{stressors.cpu: \{workers: 2, load: 100\}}, \texttt{duration: "180s"} (via \texttt{stress-ng} in target \texttt{cgroup}). & \texttt{ChaosExperiment: pod-cpu-hog}, \texttt{CPU\_\allowbreak CORES: "2"}, \texttt{CPU\_\allowbreak LOAD: "100"}, \texttt{TOTAL\_\allowbreak CHAOS\_\allowbreak DURATION: "180"} (via \texttt{go-runner} stressor). \\
\bottomrule
\end{tabularx}
\end{table*}

\subsection{Factors, Treatments, and Equivalent Fault Mapping}
We employ a \textbf{Paired Two-Factor Experimental Design} crossing:
\begin{enumerate}[leftmargin=1.2em, itemsep=2pt, topsep=3pt]
    \item \textbf{Factor A — Chaos Engineering Platform (2 Treatments):} \textbf{Treatment 1: Chaos Mesh} vs.\ \textbf{Treatment 2: LitmusChaos}.
    \item \textbf{Factor B — Stratified Fault Class (3 Strata across $N=30$ paired trials, \texttt{RANDOM\_SEED = 42}):} Following the empirical distribution of Owotogbe et al. \cite{owotogbe2026chaos}, each platform executes the exact same sequence of $30$ $(\text{target}_i, \text{fault}_i)$ pairs mapped in Table~\ref{tab:fault_equivalence} across the critical microservices $\mathcal{P}_{\text{target}} = \{\texttt{cart},\allowbreak \texttt{currency},\allowbreak \texttt{payment},\allowbreak \texttt{shipping}\}$ of the CNCF OpenTelemetry Astronomy Shop (\texttt{otel-demo}).
\end{enumerate}
To ensure strict parameter parity, both tools are configured to inject identical physical perturbations at the Linux kernel level: $2000 \pm 100\,\text{ms}$ \texttt{tc-netem} delay for $180\,\text{s}$, forced immediate pod deletion (\texttt{SIGKILL}), and 2-core $100\%$ CPU \texttt{cgroup} saturation for $180\,\text{s}$.
